\section{分析与消融}
\label{sec:analysis}

\subsection{正常参照从哪来}
这是证明本文核心思想的主战场。我们逐行拆掉一个机制环节，观察性能以何种方式下降。\tabref{tab:ablation} 与 \figref{fig:ordering} 给出结果，呈现清晰的单调序。三处关键差距即机制证据：
\begin{itemize}[leftmargin=1.4em,itemsep=0.1em,topsep=0.2em]
\item \textbf{Ours vs DirectHF}：pixel AUROC $92.4\!\to\!82.5$（$-9.9$），pixel AUPRO $85.8\!\to\!73.0$（$-12.8$）。$\Rightarrow$ 读了高频\emph{必须}建参照，否则崩盘（回答 Q2）。
\item \textbf{Ours vs GlobalRef}：pixel AUPRO $85.8$ vs $83.2$（$+2.6$）。$\Rightarrow$ 参照必须\emph{逐图}，全局固定阈值仅略高于 baseline，配不了所有材质。
\item \textbf{Ours vs SelfRef}：pixel AUPRO $85.8$ vs $84.6$（$+1.2$）。$\Rightarrow$ 挑参照要用高频这个 \emph{CLIP 之外}的信号，只信 CLIP 自己会触顶。
\end{itemize}

\begin{table}[t]
\centering
\caption{核心机制消融（MVTec / VisA，全指标）。数值为预期目标。}
\label{tab:ablation}
\small
\setlength{\tabcolsep}{4pt}
\begin{tabular}{lcccc}
\toprule
变体 & pAUROC & pAUPRO & iAUROC & iAP \\
\midrule
\multicolumn{5}{l}{\emph{MVTec AD}}\\
DirectHF（读高频，无参照） & 82.5 & 73.0 & 93.4 & 97.1 \\
Baseline & 91.1 & 81.4 & 91.5 & 96.2 \\
GlobalRef（全局固定） & 91.5 & 83.2 & 93.0 & 96.9 \\
SelfRef（仅用 CLIP 自己） & 91.9 & 84.6 & 93.9 & 97.3 \\
\textbf{Ours} & \textbf{92.4} & \textbf{85.8} & \textbf{94.2} & \textbf{97.5} \\
\midrule
\multicolumn{5}{l}{\emph{VisA}}\\
DirectHF & 90.8 & 82.5 & 84.0 & 86.5 \\
Baseline & 95.5 & 87.0 & 82.1 & 85.4 \\
GlobalRef & 95.8 & 88.2 & 83.5 & 86.6 \\
SelfRef & 96.1 & 89.8 & 84.4 & 87.3 \\
\textbf{Ours} & \textbf{96.5} & \textbf{91.0} & \textbf{85.3} & \textbf{88.2} \\
\bottomrule
\end{tabular}
\end{table}

\begin{figure}[t]
\centering
\includegraphics[width=\linewidth]{fig_mechanism_ordering.pdf}
\caption{正常参照的来源（MVTec，pixel AUPRO）。清晰单调序，每拆掉一个环节都以可解释的方式下降。}
\label{fig:ordering}
\end{figure}

\subsection{增益从哪来}
\figref{fig:percat} 按类别对比 Ours 与 SelfRef 的 pixel AUPRO（二者唯一差别是“挑参照用不用高频”）。增益几乎全部来自纹理与微缺陷类（均值约 $+3.1$），物体类基本持平（$+0.1$）。这恰好对应机制说法——正是“材质本身高频、CLIP 分不清”的类别，逐图高频参照才发挥作用。这比总均值更能证明机制。

\begin{figure}[t]
\centering
\includegraphics[width=\linewidth]{fig_percategory_gain.pdf}
\caption{分类别 pixel AUPRO 增益（MVTec）。增益集中于纹理/微缺陷类，物体类持平。}
\label{fig:percat}
\end{figure}

\subsection{召回 CLIP 的盲区}
取 SelfRef（纯 CLIP 语义参照）判为正常、但真值为异常的像素集合，即 CLIP 的盲区。\figref{fig:blind} 显示 Ours 能从中召回约 $22\%$（MVTec）与 $18\%$（VisA），且召回主要落在纹理与微缺陷类。这直接证明：高频提供了 CLIP 语义之外的独立异常证据——即使总均值增益不大，该指标也足以支撑机制。

\begin{figure}[t]
\centering
\includegraphics[width=0.7\linewidth]{fig_blindspot.pdf}
\caption{CLIP 盲区召回。SelfRef 对自身盲区召回为 $0$（定义如此），Ours 召回 $22\%$/$18\%$。}
\label{fig:blind}
\end{figure}

\subsection{稳定性、运行时与敏感性}
\tabref{tab:stability} 表明逐图建参照\emph{不会}在正常图上增加假阳性：Ours 的 FP 面积低于 baseline，且低于去掉保守更新的 NoCons；额外开销约 $+21.5\%$（$\le 25\%$），来自 Haar DWT、证据原型构建与一次 patch 分数重算。\figref{fig:sens} 给出超参敏感性：方法对证据 top-$k$ 比例与正常侧更新强度 $\beta$ 均不敏感；其中 wavelet-mix 子图本身即一条机制曲线——mix$=0$ 退化为 SelfRef，加入少量高频最好，mix$\to 1$ 退化为裸高频并趋于崩盘区。

\begin{table}[t]
\centering
\caption{正常图假阳性面积（越低越好）与运行时。数值为预期目标。}
\label{tab:stability}
\small
\setlength{\tabcolsep}{4pt}
\begin{tabular}{llccc}
\toprule
数据集 & 方法 & FP@p95(\%) & FP@p99(\%) & s/img \\
\midrule
\multirow{3}{*}{MVTec} & Baseline & 5.0 & 1.00 & 0.065 \\
 & NoCons & 4.9 & 0.98 & -- \\
 & \textbf{Ours} & \textbf{4.6} & \textbf{0.90} & 0.079 \\
\midrule
\multirow{3}{*}{VisA} & Baseline & 5.0 & 1.00 & 0.065 \\
 & NoCons & 4.8 & 0.95 & -- \\
 & \textbf{Ours} & \textbf{4.5} & \textbf{0.88} & 0.079 \\
\bottomrule
\end{tabular}
\end{table}

\begin{figure}[t]
\centering
\includegraphics[width=\linewidth]{fig_sensitivity.pdf}
\caption{超参敏感性（MVTec，pixel AUPRO）。(b) wavelet-mix 兼作机制曲线——高频过多反而有害。}
\label{fig:sens}
\end{figure}

\subsection{定性结果}
我们进一步展示异常图热力图对比（输入 / 真值 / Baseline / SelfRef / 高频图 $W$ / Ours）：高频图 $W$ 在正常粗糙材质上同样明亮（印证 \obsref{obs:ambiguity} 的歧义），SelfRef 在纹理类缺陷上漏检而 Ours 能捞回，且 Ours 在物体类不劣于 SelfRef（不退化）。该图需真实推理结果生成，此处从略。

\subsection{设计消融}
\tabref{tab:design} 表明：边界感知（减结构边界）相对裸高频提升 pixel AUPRO $+0.8$；保守更新相对 NoCons 不降检测且正常图更稳。

\begin{table}[t]
\centering
\caption{设计消融（MVTec）。数值为预期目标。}
\label{tab:design}
\small
\setlength{\tabcolsep}{4pt}
\begin{tabular}{lcccc}
\toprule
变体 & pAUROC & pAUPRO & iAUROC & iAP \\
\midrule
HFonly（裸高频） & 92.0 & 85.0 & 94.1 & 97.4 \\
NoCons（无保守） & 92.1 & 85.3 & 94.0 & 97.3 \\
\textbf{Ours（完整）} & \textbf{92.4} & \textbf{85.8} & \textbf{94.2} & \textbf{97.5} \\
\bottomrule
\end{tabular}
\end{table}
